\chapter*{Lampiran}
\addcontentsline{toc}{chapter}{Lampiran}

% ============================================================
% Lampiran A: Rubrik Penilaian
% ============================================================
\section*{Lampiran A: Bobot Penilaian dan Rubrik Komprehensif}

\subsection*{A.1 Bobot Penilaian (Sesuai RPS)}

\begin{table}[htbp]
\centering
\small
\begin{tabular}{|>{\raggedright\arraybackslash}p{3.2cm}|>{\raggedright\arraybackslash}p{5.5cm}|>{\raggedright\arraybackslash}p{3.8cm}|c|}
\hline
\textbf{Komponen} & \textbf{Teknik Asesmen} & \textbf{CPMK yang Diuji} & \textbf{Bobot (\%)} \\
\hline
Tugas/Praktikum & Lab exercise, Tugas 1--4 & Sub-CPMK 1.1--3.2 & 35 \\
\hline
Kuis (Quiz 1--4) & Quiz 1 (CPMK-1), Quiz 2 (CPMK-2), Quiz 3 (CPMK-3), Quiz 4 (CPMK-4) & CPMK-1 s.d. CPMK-4 & 10 \\
\hline
UTS & Ujian tengah (teori dan praktik, pilihan ganda) & CPMK-1, CPMK-2 & 20 \\
\hline
Proyek & Project: aplikasi web SIA sederhana & CPMK-3 (Sub-CPMK 3.1, 3.2, 3.3) & 20 \\
\hline
UAS & Ujian akhir dan presentasi (pilihan ganda) & CPMK-3, CPMK-4 & 15 \\
\hline
\textbf{Total} & & & \textbf{100} \\
\hline
\end{tabular}
\caption{Bobot Penilaian Mata Kuliah Pemrograman Web}
\end{table}

\subsection*{A.2 Rubrik Penilaian Tugas Praktik Web}

Rubrik berikut terhubung dengan CPMK dan Sub-CPMK RPS. Setiap kriteria memetakan ke capaian pembelajaran yang dapat diukur.

\begin{table}[htbp]
\centering
\small
\begin{tabular}{|>{\raggedright\arraybackslash}p{2.6cm}|>{\raggedright\arraybackslash}p{2.2cm}|>{\raggedright\arraybackslash}p{3cm}|>{\raggedright\arraybackslash}p{3cm}|>{\raggedright\arraybackslash}p{3cm}|}
\hline
\textbf{Kriteria} & \textbf{Sub-CPMK} & \textbf{Sangat Baik} & \textbf{Baik} & \textbf{Perlu Perbaikan} \\
\hline
Analisis Kebutuhan & 1.1, 1.2, 1.3 & Kebutuhan terdokumentasi jelas, wireframe/mockup profesional & Sebagian besar lengkap & Kurang lengkap \\
\hline
Struktur HTML & 2.1 & Semantik tepat, aksesibel & Semantik cukup baik & Banyak elemen tidak tepat \\
\hline
CSS dan Layout & 2.2 & Responsif, rapi, konsisten & Cukup responsif & Layout berantakan \\
\hline
Interaktivitas JS & 2.3 & Validasi dan event tepat & Ada fitur utama & Banyak error/interaksi minim \\
\hline
Library\slash Framework & 2.4 & Penggunaan tepat dan efektif & Penggunaan dasar & Tidak memadai \\
\hline
Back-end PHP & 3.1 & Logika server tepat, OOP & Sebagian fitur & Banyak error \\
\hline
Database & 3.2 & CRUD lengkap, ternormalisasi & CRUD sebagian & Banyak error \\
\hline
Integrasi API & 3.3 & Request/response terstruktur & Integrasi dasar & Tidak terintegrasi \\
\hline
Pengujian & 4.1 & Laporan lengkap, metrik terukur & Pengujian dasar & Minim pengujian \\
\hline
Keamanan & 4.2 & Validasi, prepared statement, XSS/CSRF & Validasi sebagian & Tidak ada validasi \\
\hline
Deployment & 4.3 & Konfigurasi produksi, dokumentasi & Dokumentasi dasar & Tidak ada \\
\hline
\end{tabular}
\caption{Rubrik Penilaian Tugas Praktik dan Proyek (terhubung Sub-CPMK)}
\end{table}

\subsection*{A.3 Rubrik Project SIA (CPMK-3)}

Project aplikasi web sistem informasi akademik sederhana dinilai berdasarkan:

\begin{itemize}
  \item \textbf{Sub-CPMK 3.1}: PHP OOP, session, validasi input --- back-end berfungsi dengan baik.
  \item \textbf{Sub-CPMK 3.2}: Database MySQL, CRUD, koneksi PHP--MySQL (MySQLi/PDO) --- data tersimpan dan terbaca benar.
  \item \textbf{Sub-CPMK 3.3}: Integrasi API front-end--back-end --- Fetch API mengonsumsi endpoint JSON.
\end{itemize}

Skala: A (85--100), B (70--84), C (55--69), D ($<$55). Rubrik detail mengacu tabel A.2.

% ============================================================
% Lampiran B: Contoh Template Laporan
% ============================================================
\section*{Lampiran B: Contoh Template Laporan Tugas}
\begin{enumerate}
  \item Judul dan Identitas Mahasiswa
  \item Deskripsi Masalah dan Kebutuhan Pengguna
  \item Desain UI/UX (wireframe dan mockup singkat)
  \item Implementasi (cuplikan kode penting)
  \item Hasil Pengujian dan Evaluasi
  \item Refleksi dan Kesimpulan
\end{enumerate}

% ============================================================
% Lampiran C: Glosarium
% ============================================================
\section*{Lampiran C: Glosarium Istilah Web}
\begin{itemize}
  \item \textbf{HTML}: Bahasa markup untuk struktur halaman web.
  \item \textbf{CSS}: Bahasa untuk styling dan layout halaman web.
  \item \textbf{DOM}: Representasi struktur dokumen yang dapat dimanipulasi.
  \item \textbf{API}: Antarmuka yang memungkinkan pertukaran data antar sistem.
  \item \textbf{CRUD}: Operasi Create, Read, Update, Delete pada data.
  \item \textbf{Responsive}: Desain yang menyesuaikan ukuran layar.
  \item \textbf{Session}: Mekanisme menyimpan informasi pengguna selama interaksi.
  \item \textbf{XSS/CSRF}: Jenis serangan keamanan pada aplikasi web.
\end{itemize}

% ============================================================
% Lampiran D: Referensi Tambahan
% ============================================================
\section*{Lampiran D: Referensi Tambahan}
\begin{itemize}
  \item MDN Web Docs: \url{https://developer.mozilla.org}
  \item W3Schools: \url{https://www.w3schools.com}
  \item OWASP Cheat Sheets: \url{https://cheatsheetseries.owasp.org}
\end{itemize}

